% sec1_4.tex — Section 1.4: Hypothesis Testing under Normality

Very often, the economic theory that motivated the regression equation also specifies
the values that the regression coefficients should take.  Suppose that the underlying
theory implies the restriction that $\beta_2$ equals~1.  Although Proposition~1.1 guarantees
that, on average, $b_2$ (the OLS estimate of~$\beta_2$) equals~1 if the restriction is true,
$b_2$ may not be exactly equal to~1 for a particular sample at hand.  Obviously, we cannot
conclude that the restriction is false just because the estimate~$b_2$ differs from~1.
In order for us to decide whether the sampling error $b_2 - 1$ is ``too large'' for the
restriction to be true, we need to construct from the sampling error some test statistic
whose probability distribution is known given the truth of the hypothesis.  It might
appear that doing so requires one to specify the joint distribution of
$(\vec{X}, \bm{\varepsilon})$ because, as is clear from~(1.2.14), the sampling error is a
function of $(\vec{X}, \bm{\varepsilon})$.  A surprising fact about the theory of hypothesis
testing to be presented in this section is that the distribution can be derived without
specifying the joint distribution when the conditional distribution of
$\bm{\varepsilon}$ conditional on~$\vec{X}$ is normal; there is no need to specify the
distribution of~$\vec{X}$.

In the language of hypothesis testing, the restriction to be tested (such as
``$\beta_2 = 1$'') is called the \textbf{null hypothesis} (or simply the \textbf{null}).
It is a restriction on the \textbf{maintained hypothesis}, a set of assumptions which,
combined with the null, produces some test statistic with a known distribution.  For the
present case of testing hypothesis about regression coefficients, only the normality
assumption about the conditional distribution of~$\bm{\varepsilon}$ needs to be added to the
classical regression model (Assumptions~1.1--1.4) to form the maintained hypothesis (as
just noted, there is no need to specify the joint distribution of
$(\vec{X}, \bm{\varepsilon})$).  Sometimes the maintained hypothesis is somewhat loosely
referred to as ``the model.''  We say that the model is \textbf{correctly specified} if
the maintained hypothesis is true.  Although too large a value of the test statistic is
interpreted as a failure of the null, the interpretation is valid only as long as the model
is correctly specified.  It is possible that the test statistic does not have the supposed
distribution when the null is true but the model is false.


%% ─── Normally Distributed Error Terms ─────────────────────────────────────────
\subsection*{Normally Distributed Error Terms}

In many applications, the error term consists of many miscellaneous factors not
captured by the regressors.  The Central Limit Theorem suggests that the error term has
a normal distribution.  In other applications, the error term is due to errors in
measuring the dependent variable.  It is known that very often measurement errors are
normally distributed (in fact, the normal distribution was originally developed for
measurement errors).  It is therefore worth entertaining the normality assumption:

\begin{assumption}[normality of the error term]
\label{ass:normality}
\emph{The distribution of $\bm{\varepsilon}$ conditional on $\vec{X}$ is jointly normal.}
\end{assumption}

Recall from probability theory that the normal distribution has several convenient
features:
\begin{itemize}
\item The distribution depends only on the mean and the variance.  Thus, once the mean
  and the variance are known, you can write down the density function.  If the
  distribution conditional on~$\vec{X}$ is normal, the mean and the variance can depend
  on~$\vec{X}$.  It follows that, if the distribution conditional on~$\vec{X}$ is normal
  and if neither the conditional mean nor the conditional variance depends on~$\vec{X}$,
  then the marginal (i.e., unconditional) distribution is the same normal distribution.

\item In general, if two random variables are independent, then they are uncorrelated, but
  the converse is not true.  However, if two random variables are joint normal, the
  converse is also true, so that independence and a lack of correlation are equivalent.
  This carries over to conditional distributions: if two random variables are joint
  normal and uncorrelated conditional on~$\vec{X}$, then they are independent conditional
  on~$\vec{X}$.

\item A linear function of random variables that are jointly normally distributed is
  itself normally distributed.  This also carries over to conditional distributions.  If
  the distribution of~$\bm{\varepsilon}$ conditional on~$\vec{X}$ is normal, then
  $\vec{A}\bm{\varepsilon}$, where the elements of matrix~$\vec{A}$ are functions
  of~$\vec{X}$, is normal conditional on~$\vec{X}$.
\end{itemize}

It is thanks to these features of normality that Assumption~1.5 delivers the following
properties to be exploited in the derivation of test statistics:
\begin{itemize}
\item The mean and the variance of the distribution of~$\bm{\varepsilon}$ conditional
  on~$\vec{X}$ are already specified in Assumptions~1.2 and~1.4.  Therefore,
  Assumption~1.5 together with Assumptions~1.2 and~1.4 implies that the distribution
  of~$\bm{\varepsilon}$ conditional on~$\vec{X}$ is
  $N(\vec{0}, \sigma^2 \vec{I}_n)$:
  %
  \begin{equation}
    \bm{\varepsilon} \mid \vec{X} \sim N(\vec{0}, \sigma^2 \vec{I}_n).
    \label{eq:1.4.1}
  \end{equation}
  %
  Thus, the distribution of~$\bm{\varepsilon}$ conditional on~$\vec{X}$ does not depend
  on~$\vec{X}$.  It then follows that $\bm{\varepsilon}$ and~$\vec{X}$ are
  \emph{independent}.  Therefore, in particular, the marginal or unconditional
  distribution of~$\bm{\varepsilon}$ is $N(\vec{0}, \sigma^2 \vec{I}_n)$.

\item We know from~(1.2.14) that the sampling error
  $\vec{b} - \bm{\beta}$ is linear in~$\bm{\varepsilon}$ given~$\vec{X}$.  Since
  $\bm{\varepsilon}$ is normal given~$\vec{X}$, so is the sampling error.  Its mean and
  variance are given by parts~(a) and~(b) of Proposition~1.1.  Thus, under
  Assumptions~1.1--1.5,
  %
  \begin{equation}
    (\vec{b} - \bm{\beta}) \mid \vec{X}
      \sim N(\vec{0},\, \sigma^2 \cdot (\vec{X}'\vec{X})^{-1}).
    \label{eq:1.4.2}
  \end{equation}
\end{itemize}


%% ─── Testing Hypotheses about Individual Regression Coefficients ──────────────
\subsection*{Testing Hypotheses about Individual Regression Coefficients}

The type of hypothesis we first consider is about the $k$-th coefficient
%
\[
  \mathrm{H}_0 \colon \beta_k = \overline{\beta}_k.
\]
%
Here, $\overline{\beta}_k$ is some known value specified by the null hypothesis.  We wish
to test this null against the alternative hypothesis
$\mathrm{H}_1 \colon \beta_k \neq \overline{\beta}_k$, at a significance level of~$\alpha$.
Looking at the $k$-th component of~\eqref{eq:1.4.2} and imposing the restriction of the
null, we obtain
%
\[
  (b_k - \overline{\beta}_k) \mid \vec{X}
    \sim N\!\Bigl(0,\; \sigma^2 \cdot
    \bigl((\vec{X}'\vec{X})^{-1}\bigr)_{kk}\Bigr),
\]
%
where $\bigl((\vec{X}'\vec{X})^{-1}\bigr)_{kk}$ is the $(k,k)$ element
of~$(\vec{X}'\vec{X})^{-1}$.  So if we define the ratio~$z_k$ by dividing
$b_k - \overline{\beta}_k$ by its standard deviation
%
\begin{equation}
  z_k \equiv \frac{b_k - \overline{\beta}_k}
    {\sqrt{\sigma^2 \cdot \bigl((\vec{X}'\vec{X})^{-1}\bigr)_{kk}}},
  \label{eq:1.4.3}
\end{equation}
%
then the distribution of~$z_k$ is $N(0,1)$ (the standard normal distribution).

Suppose for a second that $\sigma^2$ is known.  Then the statistic~$z_k$ has some
desirable properties as a test statistic.  First, its value can be calculated from the
sample.  Second, its distribution conditional on~$\vec{X}$ does not depend on~$\vec{X}$
(which should not be confused with the fact that the \emph{value} of~$z_k$ depends
on~$\vec{X}$).  So $z_k$ and~$\vec{X}$ are independently distributed, and, regardless of
the value of~$\vec{X}$, the distribution of~$z_k$ is the same as its unconditional
distribution.  This is convenient because different samples differ not only in~$\vec{y}$
but also in~$\vec{X}$.  Third, the distribution is known.  In particular, it does not
depend on unknown parameters (such as~$\bm{\beta}$).  (If the distribution of a statistic
depends on unknown parameters, those parameters are called \textbf{nuisance parameters}.)
Using this statistic, we can determine whether or not the sampling error
$b_k - \overline{\beta}_k$ is too large: it is too large if the test statistic takes on a
value that is surprising for a realization from the distribution.

If we do not know the true value of~$\sigma^2$, a natural idea is to replace the
nuisance parameter~$\sigma^2$ by its OLS estimate~$s^2$.  The statistic after the
substitution of~$s^2$ for~$\sigma^2$ is called the \textbf{\textit{t}-ratio} or the
\textbf{\textit{t}-value}.  The denominator of this statistic is called the
\textbf{standard error} of the OLS estimate of~$\beta_k$ and is sometimes written
as~$SE(b_k)$:
%
\begin{equation}
  SE(b_k) \equiv \sqrt{s^2 \cdot \bigl((\vec{X}'\vec{X})^{-1}\bigr)_{kk}}
    = \sqrt{(k,k) \text{ element of } \widehat{\Var(\vec{b} \mid \vec{X})}
    \text{ in (1.3.4)}}.
  \label{eq:1.4.4}
\end{equation}
%
Since $s^2$, being a function of the sample, is a random variable, this substitution
changes the distribution of the statistic, but fortunately the changed distribution, too,
is known and depends on neither nuisance parameters nor~$\vec{X}$.

\begin{proposition}[distribution of the \textit{t}-ratio]
\label{prop:t-ratio}
\emph{Suppose Assumptions~1.1--1.5 hold.  Under the null hypothesis
$\mathrm{H}_0 \colon \beta_k = \overline{\beta}_k$, the $t$-ratio defined as}
%
\begin{equation}
  t_k \equiv \frac{b_k - \overline{\beta}_k}{SE(b_k)}
    \equiv \frac{b_k - \overline{\beta}_k}
      {\sqrt{s^2 \cdot \bigl((\vec{X}'\vec{X})^{-1}\bigr)_{kk}}}
  \label{eq:1.4.5}
\end{equation}
%
\emph{is distributed as $t(n-K)$ (the $t$ distribution with $n - K$ degrees of freedom).}
\end{proposition}

\begin{proof}
We can write
%
\begin{align*}
  t_k &= \frac{b_k - \overline{\beta}_k}
    {\sqrt{\sigma^2 \cdot \bigl((\vec{X}'\vec{X})^{-1}\bigr)_{kk}}}
    \cdot \sqrt{\frac{\sigma^2}{s^2}}
    = \frac{z_k}{\sqrt{s^2/\sigma^2}} \\
  &= \frac{z_k}{\sqrt{\dfrac{\vec{e}'\vec{e}/\sigma^2}{n-K}}}
    = \frac{z_k}{\sqrt{\dfrac{q}{n-K}}},
\end{align*}
%
where $q \equiv \vec{e}'\vec{e}/\sigma^2$ to reflect the substitution of~$s^2$
for~$\sigma^2$.  We have already shown that $z_k$ is $N(0,1)$.  We will show:

\begin{enumerate}
\item[(1)] $q \mid \vec{X} \sim \chi^2(n - K)$,

\item[(2)] two random variables $z_k$ and~$q$ are independent conditional on~$\vec{X}$.
\end{enumerate}

\noindent
Then, by the definition of the $t$ distribution, the ratio of~$z_k$ to
$\sqrt{q/(n-K)}$ is distributed as~$t$ with $n - K$ degrees of freedom,%
\footnote{Fact: If $x \sim N(0,1)$, $y \sim \chi^2(m)$ and if $x$ and~$y$ are
independent, then the ratio $x/\sqrt{y/m}$ has the $t$ distribution with $m$ degrees of
freedom.}
and we are done.

\medskip
\noindent
(1) Since $\vec{e}'\vec{e} = \bm{\varepsilon}'\vec{M}\bm{\varepsilon}$ from~(1.2.12), we
have
%
\[
  q = \frac{\vec{e}'\vec{e}}{\sigma^2}
    = \frac{\bm{\varepsilon}'}{\sigma} \vec{M} \frac{\bm{\varepsilon}}{\sigma}.
\]
%
The middle matrix~$\vec{M}$, being the annihilator, is idempotent.  Also,
$\bm{\varepsilon}/\sigma \mid \vec{X} \sim N(\vec{0}, \vec{I}_n)$
by~\eqref{eq:1.4.1}.  Therefore, this quadratic form is distributed as~$\chi^2$ with
degrees of freedom equal to $\rank(\vec{M})$.%
\footnote{Fact: If $\vec{x} \sim N(\vec{0}, \vec{I}_n)$ and $\vec{A}$ is idempotent, then
$\vec{x}'\vec{A}\vec{x}$ has a chi-squared distribution with degrees of freedom equal to
the rank of~$\vec{A}$.}
But $\rank(\vec{M}) = \tr(\vec{M})$, because $\vec{M}$ is idempotent.%
\footnote{Fact: If $\vec{A}$ is idempotent, then $\rank(\vec{A}) = \tr(\vec{A})$.}
We have already shown in the proof of Proposition~1.2 that
$\tr(\vec{M}) = n - K$.  So $q \mid \vec{X} \sim \chi^2(n - K)$.

\medskip
\noindent
(2) Both $\vec{b}$ and~$\vec{e}$ are linear functions of~$\bm{\varepsilon}$
(by~(1.2.14) and the fact that $\vec{e} = \vec{M}\bm{\varepsilon}$), so they are jointly
normal conditional on~$\vec{X}$.  Also, they are uncorrelated conditional on~$\vec{X}$
(see part~(d) of Proposition~1.1).  So $\vec{b}$ and~$\vec{e}$ are independently
distributed conditional on~$\vec{X}$.  But $z_k$ is a function of~$\vec{b}$ and $q$ is a
function of~$\vec{e}$.  So $z_k$ and~$q$ are independently distributed conditional
on~$\vec{X}$.%
\footnote{Fact: If $\vec{x}$ and~$\vec{y}$ are independently distributed, then so are
$f(\vec{x})$ and $g(\vec{y})$.}
\end{proof}


%% ─── Decision Rule for the t-Test ─────────────────────────────────────────────
\subsection*{Decision Rule for the \textit{t}-Test}

The test of the null hypothesis based on the $t$-ratio is called the $t$-test and proceeds
as follows:

\begin{description}[style=unboxed, leftmargin=0pt]
\item[\textit{Step~1:}] Given the hypothesized value, $\overline{\beta}_k$, of~$\beta_k$,
  form the $t$-ratio as in~\eqref{eq:1.4.5}.  Too large a deviation of~$t_k$ from~0 is a
  sign of the failure of the null hypothesis.  The next step specifies how large is too
  large.

\item[\textit{Step~2:}] Go to the $t$-table (most statistics and econometrics textbooks
  include the $t$-table) and look up the entry for $n - K$ degrees of freedom.  Find the
  \textbf{critical value}, $t_{\alpha/2}(n - K)$, such that the area in the $t$
  distribution to the right of $t_{\alpha/2}(n - K)$ is~$\alpha/2$, as illustrated in
  Figure~\ref{fig:1.3}.  (If $n - K = 30$ and $\alpha = 5\%$, for example,
  $t_{\alpha/2}(n - K) = 2.042$.)  Then, since the $t$ distribution is symmetric around~0,
  %
  \[
    \mathrm{Prob}\bigl(-t_{\alpha/2}(n-K) < t < t_{\alpha/2}(n-K)\bigr) = 1 - \alpha.
  \]

\item[\textit{Step~3:}] Accept $\mathrm{H}_0$ if
  $-t_{\alpha/2}(n-K) < t_k < t_{\alpha/2}(n-K)$ (that is, if
  $|t_k| < t_{\alpha/2}(n-K)$), where $t_k$ is the $t$-ratio from \emph{Step~1}.  Reject
  $\mathrm{H}_0$ otherwise.  Since $t_k \sim t(n-K)$ under~$\mathrm{H}_0$, the
  probability of rejecting~$\mathrm{H}_0$ when $\mathrm{H}_0$ is true is~$\alpha$.  So
  the size (significance level) of the test is indeed~$\alpha$.
\end{description}

A convenient feature of the $t$-test is that the critical value does not depend
on~$\vec{X}$; there is no need to calculate critical values for each sample.

\begin{figure}[htbp]
\centering
\includegraphics[width=0.9\textwidth,trim=50 400 50 400,clip]{figures/fig_1_3.png}
\caption{$t$ Distribution}
\label{fig:1.3}
\end{figure}


%% ─── Confidence Interval ──────────────────────────────────────────────────────
\subsection*{Confidence Interval}

\emph{Step~3} can also be stated in terms of $b_k$ and~$SE(b_k)$.  Since $t_k$ is as
in~\eqref{eq:1.4.5}, you accept~$\mathrm{H}_0$ whenever
%
\[
  -t_{\alpha/2}(n-K) < \frac{b_k - \overline{\beta}_k}{SE(b_k)} < t_{\alpha/2}(n-K)
\]
%
or
%
\[
  b_k - SE(b_k) \cdot t_{\alpha/2}(n-K)
    < \overline{\beta}_k
    < b_k + SE(b_k) \cdot t_{\alpha/2}(n-K).
\]
%
Therefore, we accept if and only if the hypothesized value~$\overline{\beta}_k$ falls in
the interval:
%
\begin{equation}
  \bigl[\, b_k - SE(b_k) \cdot t_{\alpha/2}(n-K),\;
    b_k + SE(b_k) \cdot t_{\alpha/2}(n-K) \,\bigr].
  \label{eq:1.4.6}
\end{equation}
%
This interval is called the \textbf{level $\boldsymbol{1 - \alpha}$ confidence interval}.
It is narrower the smaller the standard error.  Thus, the smallness of the standard error
is a measure of the estimator's precision.


%% ─── p-Value ──────────────────────────────────────────────────────────────────
\subsection*{\textit{p}-Value}

The decision rule of the $t$-test can also be stated using the \textbf{\textit{p}-value}.

\begin{description}[style=unboxed, leftmargin=0pt]
\item[\textit{Step~1:}] Same as above.

\item[\textit{Step~2:}] Rather than finding the critical value $t_{\alpha/2}(n - K)$,
  calculate
  %
  \[
    p = \mathrm{Prob}(t > |t_k|) \times 2.
  \]
  %
  Since the $t$ distribution is symmetric around~0,
  $\mathrm{Prob}(t > |t_k|) = \mathrm{Prob}(t < -|t_k|)$, so
  %
  \begin{equation}
    \mathrm{Prob}(-|t_k| < t < |t_k|) = 1 - p.
    \label{eq:1.4.7}
  \end{equation}

\item[\textit{Step~3:}] Accept $\mathrm{H}_0$ if $p > \alpha$.  Reject otherwise.
\end{description}

To see the equivalence of the two decision rules, one based on the critical values such
as~$t_{\alpha/2}(n - K)$ and the other based on the $p$-value, refer to
Figure~\ref{fig:1.3}.  If $\mathrm{Prob}(t > |t_k|)$ is greater than~$\alpha/2$ (as in the
figure), that is, if the $p$-value is more than~$\alpha$, then $|t_k|$ must be to the left
of~$t_{\alpha/2}(n - K)$.  This means from \emph{Step~3} that the null hypothesis is not
rejected.  Thus, when $p$ is small, the $t$-ratio is surprisingly large for a random
variable from the $t$ distribution.  The smaller the~$p$, the stronger the rejection.

Examples of the $t$-test can be found in Section~1.7.


%% ─── Linear Hypotheses ────────────────────────────────────────────────────────
\subsection*{Linear Hypotheses}

The null hypothesis we wish to test may not be a restriction about individual regression
coefficients of the maintained hypothesis; it is often about linear combinations of them
written as a system of linear equations:
%
\begin{equation}
  \mathrm{H}_0 \colon \vec{R}\bm{\beta} = \vec{r},
  \label{eq:1.4.8}
\end{equation}
%
where values of~$\vec{R}$ and~$\vec{r}$ are known and specified by the hypothesis.  We
denote the number of equations, which is the dimension of~$\vec{r}$, by~$\#\vec{r}$.  So
$\vec{R}$ is $\#\vec{r} \times K$.  These $\#\vec{r}$ equations are restrictions on the
coefficients in the maintained hypothesis.  It is called a linear hypothesis because each
equation is linear.  To make sure that there are no redundant equations and that the
equations are consistent with each other, we require that
$\rank(\vec{R}) = \#\vec{r}$ (i.e., $\vec{R}$ is of full \emph{row} rank with its rank
equaling the number of rows).  But do not be too conscious about the rank condition; in
specific applications, it is very easy to spot a failure of the rank condition if there is
one.

\begin{example}[continuation of Example~1.2]
\label{ex:linear-hyp}
Consider the wage equation of Example~1.2 where $K = 4$.  We might wish to test the
hypothesis that education and tenure have equal impact on the wage rate and that there is
no experience effect.  The hypothesis is two equations (so $\#\vec{r} = 2$):
%
\[
  \beta_2 = \beta_3 \quad \text{and} \quad \beta_4 = 0.
\]
%
This can be cast in the format $\vec{R}\bm{\beta} = \vec{r}$ if $\vec{R}$ and~$\vec{r}$
are defined as
%
\[
  \vec{R} = \begin{bmatrix} 0 & 1 & -1 & 0 \\ 0 & 0 & 0 & 1 \end{bmatrix}, \quad
  \vec{r} = \begin{bmatrix} 0 \\ 0 \end{bmatrix}.
\]
%
Because the two rows of this~$\vec{R}$ are linearly independent, the rank condition is
satisfied.

But suppose we require additionally that
%
\[
  \beta_2 - \beta_3 = \beta_4.
\]
%
This is redundant because it holds whenever the first two equations do.  With these three
equations, $\#\vec{r} = 3$ and
%
\[
  \vec{R} = \begin{bmatrix} 0 & 1 & -1 & 0 \\ 0 & 0 & 0 & 1 \\ 0 & 1 & -1 & -1
  \end{bmatrix}, \quad
  \vec{r} = \begin{bmatrix} 0 \\ 0 \\ 0 \end{bmatrix}.
\]
%
Since the third row of~$\vec{R}$ is the difference between the first two, $\vec{R}$ is
not of full row rank.  The consequence of adding redundant equations is that $\vec{R}$ no
longer meets the full row rank condition.

As an example of inconsistent equations, consider adding to the first two equations the
third equation $\beta_4 = 0.5$.  Evidently, $\beta_4$ cannot be~0 and~0.5 at the same
time.  The hypothesis is inconsistent because there is no~$\bm{\beta}$ that satisfies the
three equations simultaneously.  If we nevertheless included this equation, then $\vec{R}$
and~$\vec{r}$ would become
%
\[
  \vec{R} = \begin{bmatrix} 0 & 1 & -1 & 0 \\ 0 & 0 & 0 & 1 \\ 0 & 0 & 0 & 1
  \end{bmatrix}, \quad
  \vec{r} = \begin{bmatrix} 0 \\ 0 \\ 0.5 \end{bmatrix}.
\]
%
Again, the full row rank condition is not satisfied because the rank of~$\vec{R}$ is~2
while $\#\vec{r} = 3$.
\end{example}


%% ─── The F-Test ───────────────────────────────────────────────────────────────
\subsection*{The \textit{F}-Test}

To test linear hypotheses, we look for a test statistic that has a known distribution
under the null hypothesis.

\begin{proposition}[distribution of the \textit{F}-ratio]
\label{prop:F-ratio}
\emph{Suppose Assumptions~1.1--1.5 hold.  Under the null hypothesis
$\mathrm{H}_0 \colon \vec{R}\bm{\beta} = \vec{r}$, where $\vec{R}$ is
$\#\vec{r} \times K$ with $\rank(\vec{R}) = \#\vec{r}$, the \textbf{\textit{F}-ratio}
defined as}
%
\begin{align}
  F &\equiv \frac{(\vec{R}\vec{b} - \vec{r})'
    \bigl[\vec{R}(\vec{X}'\vec{X})^{-1}\vec{R}'\bigr]^{-1}
    (\vec{R}\vec{b} - \vec{r}) / \#\vec{r}}{s^2} \notag \\
    &= (\vec{R}\vec{b} - \vec{r})'
    [\vec{R}\widehat{\Var(\vec{b} \mid \vec{X})}\vec{R}']^{-1}
    (\vec{R}\vec{b} - \vec{r}) / \#\vec{r}
    \quad \text{(by (1.3.4))}
  \label{eq:1.4.9}
\end{align}
%
\emph{is distributed as $F(\#\vec{r}, n - K)$ (the $F$ distribution with $\#\vec{r}$
and $n - K$ degrees of freedom).}
\end{proposition}

As in Proposition~1.3, it suffices to show that the distribution conditional on~$\vec{X}$
is $F(\#\vec{r}, n - K)$; because the $F$ distribution does not depend on~$\vec{X}$, it
is also the unconditional distribution of the statistic.

\begin{proof}
Since $s^2 = \vec{e}'\vec{e}/(n - K)$, we can write
%
\[
  F = \frac{w / \#\vec{r}}{q / (n - K)}
\]
%
where
%
\[
  w \equiv (\vec{R}\vec{b} - \vec{r})'
    [\sigma^2 \cdot \vec{R}(\vec{X}'\vec{X})^{-1}\vec{R}']^{-1}
    (\vec{R}\vec{b} - \vec{r})
  \quad \text{and} \quad
  q \equiv \frac{\vec{e}'\vec{e}}{\sigma^2}.
\]
%
We need to show
%
\begin{enumerate}
\item[(1)] $w \mid \vec{X} \sim \chi^2(\#\vec{r})$,

\item[(2)] $q \mid \vec{X} \sim \chi^2(n - K)$ (this is part~(1) in the proof of
  Proposition~1.3),

\item[(3)] $w$ and~$q$ are independently distributed conditional on~$\vec{X}$.
\end{enumerate}

\noindent
Then, by the definition of the $F$ distribution, the $F$-ratio
$\sim F(\#\vec{r}, n - K)$.%
\footnote{Fact: Let $\vec{x}$ be an $m$ dimensional random vector.  If
$\vec{x} \sim N(\bm{\mu}, \bm{\Sigma})$ with $\bm{\Sigma}$ nonsingular, then
$(\vec{x} - \bm{\mu})'\bm{\Sigma}^{-1}(\vec{x} - \bm{\mu}) \sim \chi^2(m)$.}

\medskip
\noindent
(1) Let $\vec{v} \equiv \vec{R}\vec{b} - \vec{r}$.  Under~$\mathrm{H}_0$,
$\vec{R}\vec{b} - \vec{r} = \vec{R}(\vec{b} - \bm{\beta})$.  So by~\eqref{eq:1.4.2},
conditional on~$\vec{X}$, $\vec{v}$ is normal with mean~$\vec{0}$, and its variance is
given by
%
\[
  \Var(\vec{v} \mid \vec{X})
    = \Var(\vec{R}(\vec{b} - \bm{\beta}) \mid \vec{X})
    = \vec{R}\,\Var(\vec{b} - \bm{\beta} \mid \vec{X})\,\vec{R}'
    = \sigma^2 \cdot \vec{R}(\vec{X}'\vec{X})^{-1}\vec{R}',
\]
%
which is none other than the inverse of the middle matrix in the quadratic form for~$w$.
Hence, $w$ can be written as
$\vec{v}'\,\Var(\vec{v} \mid \vec{X})^{-1}\,\vec{v}$.  Since $\vec{R}$ is of full row
rank and $\vec{X}'\vec{X}$ is nonsingular,
$\sigma^2 \cdot \vec{R}(\vec{X}'\vec{X})^{-1}\vec{R}'$ is nonsingular (why?  Showing
this is a review question).  Therefore, by the definition of the $\chi^2$ distribution,
$w \mid \vec{X} \sim \chi^2(\#\vec{r})$.

\medskip
\noindent
(3) $w$ is a function of~$\vec{b}$ and $q$ is a function of~$\vec{e}$.  But $\vec{b}$
and~$\vec{e}$ are independently distributed conditional on~$\vec{X}$, as shown in
part~(2) of the proof of Proposition~1.3.  So $w$ and~$q$ are independently distributed
conditional on~$\vec{X}$.
\end{proof}

If the null hypothesis $\vec{R}\bm{\beta} = \vec{r}$ is true, we expect
$\vec{R}\vec{b} - \vec{r}$ to be small, so large values of~$F$ should be taken as
evidence for a failure of the null.  This means that we look at only the upper tail of the
distribution in the $F$-statistic.  The decision rule of the $F$-test at the significance
level of~$\alpha$ is as follows.

\begin{description}[style=unboxed, leftmargin=0pt]
\item[\textit{Step~1:}] Calculate the $F$-ratio by the formula~\eqref{eq:1.4.9}.

\item[\textit{Step~2:}] Go to the table of $F$ distribution and look up the entry for
  $\#\vec{r}$ (the numerator degrees of freedom) and $n - K$ (the denominator degrees of
  freedom).  Find the critical value $F_\alpha(\#\vec{r}, n - K)$ that leaves~$\alpha$ for
  the upper tail of the $F$ distribution, as illustrated in Figure~\ref{fig:1.4}.  For
  example, when $\#\vec{r} = 3$, $n - K = 30$, and $\alpha = 5\%$, the critical value
  $F_{.05}(3, 30)$ is~2.92.

\item[\textit{Step~3:}] Accept the null if the $F$-ratio from \emph{Step~1} is less than
  $F_\alpha(\#\vec{r}, n - K)$.  Reject otherwise.
\end{description}

This decision rule can also be described in terms of the $p$-value:

\begin{description}[style=unboxed, leftmargin=0pt]
\item[\textit{Step~1:}] Same as above.

\item[\textit{Step~2:}] Calculate
  %
  \[
    p = \text{area of the upper tail of the $F$ distribution to the right of the
    $F$-ratio.}
  \]

\item[\textit{Step~3:}] Accept the null if $p > \alpha$; reject otherwise.
\end{description}

Thus, a \emph{small} $p$-value is a signal of the failure of the null.

\begin{figure}[htbp]
\centering
\includegraphics[width=0.9\textwidth,trim=50 400 50 400,clip]{figures/fig_1_4.png}
\caption{$F$ Distribution}
\label{fig:1.4}
\end{figure}


%% ─── A More Convenient Expression for F ───────────────────────────────────────
\subsection*{A More Convenient Expression for \textit{F}}

The above derivation of the $F$-ratio is by the \textbf{Wald principle}, because it is
based on the unrestricted estimator, which is not constrained to satisfy the restrictions
of the null hypothesis.  Calculating the $F$-ratio by the formula~\eqref{eq:1.4.9}
requires matrix inversion and multiplication.  Fortunately, there is a convenient
alternative formula involving two different sum of squared residuals: one is \emph{SSR},
the minimized sum of squared residuals from~(1.2.1) now denoted as~$SSR_U$, and the other
is the restricted sum of squared residuals, denoted~$SSR_R$, obtained from
%
\begin{equation}
  \min_{\widetilde{\bm{\beta}}} SSR(\widetilde{\bm{\beta}})
    \quad \text{s.t.} \quad \vec{R}\widetilde{\bm{\beta}} = \vec{r}.
  \label{eq:1.4.10}
\end{equation}
%
Finding the $\widetilde{\bm{\beta}}$ that achieves this constrained minimization is called
the \textbf{restricted regression} or \textbf{restricted least squares}.  It is left as an
analytical exercise to show that the $F$-ratio equals
%
\begin{equation}
  F = \frac{(SSR_R - SSR_U) / \#\vec{r}}{SSR_U / (n - K)},
  \label{eq:1.4.11}
\end{equation}
%
which is the difference in the objective function deflated by the estimate of the error
variance.  This derivation of the $F$-ratio is analogous to how the likelihood-ratio
statistic is derived in maximum likelihood estimation as the difference in log likelihood
with and without the imposition of the null hypothesis.  For this reason, this second
derivation of the $F$-ratio is said to be by the \textbf{Likelihood-Ratio principle}.
There is a closed-form expression for the restricted least squares estimator
of~$\bm{\beta}$.  Deriving the expression is left as an analytical exercise.  The
computation of restricted least squares will be explained in the context of the empirical
example in Section~1.7.


%% ─── t versus F ───────────────────────────────────────────────────────────────
\subsection*{\textit{t} versus \textit{F}}

Because hypotheses about individual coefficients are linear hypotheses, the $t$-test of
$\mathrm{H}_0 \colon \beta_k = \overline{\beta}_k$ is a special case of the $F$-test.  To
see this, note that the hypothesis can be written as $\vec{R}\bm{\beta} = \vec{r}$ with
%
\[
  \underset{(1 \times K)}{\vec{R}}
    = \bigl[\, 0 \;\; \cdots \;\; 0 \;\; \underset{(k)}{1} \;\; 0 \;\;
    \cdots \;\; 0 \,\bigr], \quad
  \vec{r} = \overline{\beta}_k.
\]
%
So by~\eqref{eq:1.4.9} the $F$-ratio is
%
\[
  F = (b_k - \overline{\beta}_k)
    \bigl[s^2 \cdot (k,k) \text{ element of } (\vec{X}'\vec{X})^{-1}\bigr]^{-1}
    (b_k - \overline{\beta}_k),
\]
%
which is the square of the $t$-ratio in~\eqref{eq:1.4.5}.  Since a random variable
distributed as $F(1, n - K)$ is the square of a random variable distributed as~$t(n-K)$,
the $t$- and $F$-tests give the same test result.

Sometimes, the null is that a set of individual regression coefficients equal certain
values.  For example, assume $K = 2$ and consider
%
\[
  \mathrm{H}_0 \colon \beta_1 = 1 \quad \text{and} \quad \beta_2 = 0.
\]
%
This can be written as a linear hypothesis $\vec{R}\bm{\beta} = \vec{r}$ for
$\vec{R} = \vec{I}_2$ and $\vec{r} = (1, 0)'$.  So the $F$-test can be used.  It is
tempting, however, to conduct the $t$-test separately for each individual coefficient of
the hypothesis.  We might accept~$\mathrm{H}_0$ if both restrictions $\beta_1 = 1$ and
$\beta_2 = 0$ pass the $t$-test.  This amounts to using the confidence region of
%
\[
  \bigl\{(\beta_1, \beta_2) \mid
    b_1 - SE(b_1) \cdot t_{\alpha/2}(n-K) < \beta_1
    < b_1 + SE(b_1) \cdot t_{\alpha/2}(n-K),
\]
\[
  \qquad b_2 - SE(b_2) \cdot t_{\alpha/2}(n-K) < \beta_2
    < b_2 + SE(b_2) \cdot t_{\alpha/2}(n-K) \bigr\},
\]
%
which is a rectangular region in the $(\beta_1, \beta_2)$ plane, as illustrated in
Figure~\ref{fig:1.5}.  If $(1, 0)$, the point in the $(\beta_1, \beta_2)$ plane specified
by the null, falls in this region, one would accept the null.  On the other hand, the
confidence region for the $F$-test is
%
\[
  \Bigl\{(\beta_1, \beta_2) \mid
    (b_1 - \beta_1,\; b_2 - \beta_2)\,
    \widehat{\Var(\vec{b} \mid \vec{X})}^{-1}
    \begin{bmatrix} b_1 - \beta_1 \\ b_2 - \beta_2 \end{bmatrix}
    < 2 F_\alpha(\#\vec{r}, n - K) \Bigr\}.
\]
%
Since $\widehat{\Var(\vec{b} \mid \vec{X})}$ is positive definite, the $F$-test acceptance
region is an ellipse in the $(\beta_1, \beta_2)$ plane.  The two confidence regions look
typically like Figure~\ref{fig:1.5}.

The $F$-test should be preferred to the test using two $t$-ratios for two reasons.
First, if the size (significance level) in each of the two $t$-tests is~$\alpha$, then
the overall size (the probability that $(1, 0)$ is outside the rectangular region) is
not~$\alpha$.  Second, as will be noted in the next section (see~(1.5.19)), the $F$-test
is a likelihood ratio test and likelihood-ratio tests have certain desirable properties.
So even if the significance level in each $t$-test is controlled so that the overall size
is~$\alpha$, the test is less desirable than the $F$-test.%
\footnote{For more details on the relationship between the $t$-test and the $F$-tests,
see Scheff\'{e} (1959, p.~46).}

\begin{figure}[htbp]
\centering
\includegraphics[width=0.9\textwidth,trim=50 400 50 400,clip]{figures/fig_1_5.png}
\caption{$t$- versus $F$-Tests}
\label{fig:1.5}
\end{figure}


%% ─── An Example of a Test Statistic Whose Distribution Depends on X ───────────
\subsection*{An Example of a Test Statistic Whose Distribution Depends on X}

To place the discussion of this section in a proper perspective, it may be useful to note
that there are some statistics whose conditional distribution depends on~$\vec{X}$.
Consider the celebrated \textbf{Durbin--Watson statistic}:
%
\[
  \frac{\sum_{i=2}^{n}(e_i - e_{i-1})^2}{\sum_{i=1}^{n} e_i^2}.
\]
%
The conditional distribution, and hence the critical values, of this statistic depend
on~$\vec{X}$, but J.~Durbin and G.\,S.~Watson have shown that the critical values fall
between two bounds (which depends on the sample size, the number of regressors, and
whether the regressor includes a constant).  Therefore, the critical values for the
unconditional distribution, too, fall between these bounds.

The statistic is designed for testing whether there is no serial correlation in the error
term.  Thus, the null hypothesis is Assumption~1.4, while the maintained hypothesis is the
other assumptions of the classical regression model (including the strict exogeneity
assumption) and the normality assumption.  But, as emphasized in Section~1.1, the strict
exogeneity assumption is not satisfied in time-series models typically encountered in
econometrics, and serial correlation is an issue that arises only in time-series models.
Thus, the Durbin--Watson statistic is not useful in econometrics.  More useful tests for
serial correlation, which are all based on large-sample theory, will be covered in the
next chapter.
